\scanpage{71}
\mspage{67}
% source lines 1--8: cancelled draft, with horizontal and diagonal strokes
\deleted{\unsure{Ce|ce} qui signifie \unsure{d}}

\deleted{\uline{Définition}, Un couple $(\alpha_*,\beta_*)$ d'applications}
\deleted{\inserted{application}}
\diagram{\begin{tikzcd}[ampersand replacement=\&]
\deleted{\alpha_*:\Sigma_E} \arrow[r] \& \deleted{\Sigma_{E'}}
\end{tikzcd}}
\deleted{est dite une \uline{correspondance} \uline{\unsure{disjonction|disjonctive}} de $E$ à $E'$, si \deleted{il}\inserted{elle} \deleted{existe une application}\inserted{\unsure{commute aux}}}
\[
\deleted{\beta_*:\Sigma_{E'}\unsure{:}}
\]
\deleted{$\operatorname{Sup}$ quelconques, et s'il existe une application}

% source lines 9--15
\uline{\unsure{Suite|Sottise}} Une correspondance $\alpha$ entre deux ensembles $E,E'$ (i.e. une \uline{relation} entre $E,E'$ i.e. une partie $\alpha\subset E\times E'$) peut être interprétée comme une application
\diagram{\begin{tikzcd}[ampersand replacement=\&]
\alpha_*:\unsure{\mathcal{P}|\mathscr{P}}(E) \arrow[r] \& \mathcal{P}(E')
\end{tikzcd}}
qui commute aux $\bigcup_i$ quelconques.
% margin note written slanted
\margin{Composition des corr. $\Leftrightarrow$ \unsure{composition} des applications \unsure{sur|les} $\mathcal{P}(\unsure{E|E'})$.}

% source lines 16--22
De même, \deleted{\unsure{une}} si $E,E'$ sont des disjoncteurs, je vais considérer les applications
\diagram{\begin{tikzcd}[ampersand replacement=\&]
\alpha_*:\Sigma_E \arrow[r,no head] \& \unsure{\Sigma}_{E'}
\end{tikzcd}}
qui commutent aux $\operatorname{Sup}$ quelconques, et qui généralisent donc les correspondances de disjoncteurs \unsure{``}triviaux''. Il suffit de se donner l'application
\diagram{\begin{tikzcd}[ampersand replacement=\&]
\alpha:E \arrow[r] \& \Sigma_{E'}\\
x \arrow[r,maps to] \& \alpha_*(\underbrace{\operatorname{supp}^{\circ}_E(\{x\})}_{\unsure{S|s}(x)})
\end{tikzcd}}
% archive sheet number: 70
